%!TEX root = ../AdvMathLecture.tex
%
% Section 10.1: Line Integrals
% 

%%~~~~~~~~~~~~~~~~~~~~~~   []   ~~~~~~~~~~~~~~~~~~~~~~~~~~~
%\begin{frame}
%\frametitle{\Large {\color[rgb]{0.02,0.20,1.0} \UL{Parametric Curves \& Integration Review}}}
%\scriptsize
%
%\begin{block}{Key Concepts:}
%%-------LIST---------
%\begin{itemize}
%	\item[...] Parametric Curves
%	\item[...] Review: Integration
%	\item[...] Arclength of a Curve
%\end{itemize}
%%-------------------
%\end{block}
%
%\vspace{0.6cm}
%
%\begin{block}{Mathematical Tools:}
%%-------LIST---------
%\begin{itemize}
%	\item[...] Parametrizing curves
%	\item[...] Compute the tangent vector of a curve
%	\item[...] Deriving and Interpreting Integrals
%	\item[...] Compute the arclength of a curve
%\end{itemize}
%%-------------------
%\end{block}
%
%%\vspace{3.0cm}
%\end{frame}

%%~~~~~~~~~~~~~~~~~~~~~~~~~~~~~~~~~~~~~~~~~~~~~~~~~~~~~~~~~
%\begin{frame}
%\frametitle{}
%\scriptsize
%\begin{columns}
%\begin{column}{12.0cm}
%	Derive the Integral Expression for:
%	
%	\OBUL Disc/Washer Method
%	
%	\OBUL Shell Method
%	
%	\OBUL Work/Energy
%	
%	\OBUL Mass
%	
%	\OBUL Surface Area
%	
%	\OBUL Volume
%	
%	\OBUL Arc Length
%	
%	\OBUL Position
%	
%	\OBUL Velocity
%	
%\end{column}
%\begin{column}{0.0cm}\end{column}
%\end{columns}
%\end{frame}

%~~~~~~~~~~~~~~~~~~~~~~~~~~~~~~~~~~~~~~~~~~~~~~~~~~~~~~~~~
\begin{frame}
\frametitle{Parametric Curves \& Integration Review}
\begin{columns}
\begin{column}{12.0cm}

{I. Parametric Curves}

\vspace{0.2cm}

\scriptsize

	\DEF A {\bf parametric curve} is a curve in space that depends on a single variable which we called the parameter.  In the figure below, $t$ is the parameter.
	
\vspace{0.3cm}

\begin{columns}
\begin{column}{3.0cm}

\begin{center}
	\includegraphics[scale = 0.25]{Figures/9_5_Fig1.png} \\
\end{center}
		
\end{column}
\begin{column}{6.5cm}

\vspace{0.2cm}

${\bf r}(t) = [x(t),	y(t),	z (t)] = x(t)\VI + y(t)\VJ + z(t)\VK$

\vspace{0.2cm}

\OBUL If $t$ is defined on a finite interval $[a,b]$ then \NEWLINE $C$ will be a finite curve.

\vspace{0.2cm}

\OBUL If $t$ is defined on a infinite interval, say \NEWLINE $[0,\infty]$, then $C$ will be a infinite or closed curve.

\begin{center}
	\includegraphics[scale = 0.14]{Figures/9_5_Fig4.png} \quad	\includegraphics[scale = 0.12]{Figures/9_5_Fig3.png} \\
	finite curve \hspace{1.5cm} closed curve
\end{center}
	
\end{column}
\end{columns}

\vspace{0.6cm}
	
\UL{Note}: Defining a curve parametrically gives the curve an orientation (direction). \\

\vspace{0.2cm}
	
\end{column}
\begin{column}{0.0cm}\end{column}
\end{columns}
\end{frame}

%~~~~~~~~~~~~~~~~~~~~~~~~~~~~~~~~~~~~~~~~~~~~~~~~~~~~~~~~~
\begin{frame}
{}
\scriptsize
\begin{columns}
\begin{column}{12.0cm}
	Find {\bf r}(t) by parametrizing the following curves
	
	\vspace{0.2cm}
	
	\EX $x = y^2 + y - 15,\,\, z = 0$ \quad (Intersection of 
	\ON<2->{\B{\,\, a parabolic cylinder and plane  \,\,}})
	
	\vspace{0.2cm}
	
	The easy parameterization is to let \, 
	\ON<3->{\B{$t=y$:}}
	
	\vspace{0.2cm}
	
	\begin{columns}
	\begin{column}{1.0cm}
		%--------TABLE---------
		\begin{tabular}{l}
			$x = \ON<5->{\B{t^2 + t - 15}}$ \\
			$y = \ON<4->{\B{t}}$ \\
			$z = \ON<6->{\B{0}}$
		\end{tabular}
	\end{column}
	\begin{column}{0.3cm}
		$\D{\implies}$
	\end{column}
	\begin{column}{6.3cm}
		${\bf r}(t) = 
		\ON<7->{\B{[\,t^2 + t - 15\,,\,t\,,\,0\,]}}$ \qquad $t \in 
		\ON<8->{\B{\RE}}$
	\end{column}
	\end{columns}
	
	\vspace{0.9cm}
	
	\EX $4x - y = -1,\,\, 5x - z = 0$ \quad (Intersection of 
	\ON<9->{\B{\,\, 2 planes \,\,}})
	
	\vspace{0.2cm}
	
	Let \, 
	\ON<10->{\B{$t=x$}} 
	\ON<11->{\B{\, and solve $y$ \& $z$ in terms of $x$ :}}
	
	\vspace{0.2cm}
	
	\begin{columns}
	\begin{column}{1.0cm}
		%--------TABLE---------
		\begin{tabular}{l}
			$x = \ON<12->{\B{t}}$ \\
			$y = \ON<13->{\B{4t+1}}$ \\
			$z = \ON<14->{\B{5t}}$
		\end{tabular}
	\end{column}
	\begin{column}{0.3cm}
		$\D{\implies}$
	\end{column}
	\begin{column}{6.3cm}
		${\bf r}(t) = 
		\ON<15->{\B{[\,t\,,\,4t+1\,,\,5t\,]}}$ \qquad $t \in 
		\ON<16->{\B{\RE}}$
	\end{column}
	\end{columns}
	
	\vspace{0.9cm}
	
	\EX $x^2 + y^2 = 2$ , $z = -2$ \quad (Intersection of 
	\ON<17->{\B{\,\, a cylinder and plane \,\,}} 
	\ON<18->{\B{- \,\, circle of radius $\sqrt{2}$ \@ $z=-2$ \,\,}})
	
	\vspace{0.2cm}
	
	Let \,\ON<19->{\B{$t$ be the polar angle in the $xy$-plane:}} \,
	
	\vspace{0.2cm}
	
	\begin{columns}
	\begin{column}{1.0cm}
		%--------TABLE---------
		\begin{tabular}{l}
			$x = \ON<20->{\B{\sqrt{2}\cos t}}$ \\
			$y = \ON<21->{\B{\sqrt{2}\sin t}}$ \\
			$z = \ON<22->{\B{-2}}$
		\end{tabular}
	\end{column}
	\begin{column}{0.3cm}
		$\D{\implies}$
	\end{column}
	\begin{column}{6.7cm}
		${\bf r}(t) = 
		\ON<23->{\B{[\,\sqrt{2}\cos t\,,\,\sqrt{2}\sin t\,,\,-2\,]}}$ \qquad $t \in 
		\ON<24->{\B{[0,2\pi]}}$
	\end{column}
	\end{columns}
	
	\vspace{0.9cm}
	
\end{column}
\begin{column}{0.0cm}\end{column}
\end{columns}
\end{frame}

%~~~~~~~~~~~~~~~~~~~~~~~~~~~~~~~~~~~~~~~~~~~~~~~~~~~~~~~~~
\begin{frame}
{}
\scriptsize
\begin{columns}
\begin{column}{12.0cm}
	
	\vspace{0.2cm}
	
	\EX $x = y^3 + y,\,\, z = x - y$ \quad (Intersection of 
	\ON<2->{\B{\,\, a cubic cylinder and plane  \,\,}})
	
	\vspace{0.2cm}
	
	The easy parameterization is to let \, 
	\ON<3->{\B{$t=y$:}}
	
	\vspace{0.2cm}
	
	\begin{columns}
	\begin{column}{1.0cm}
		%--------TABLE---------
		\begin{tabular}{l}
			$x = \ON<5->{\B{t^3 + t}}$ \\
			$y = \ON<4->{\B{t}}$ \\
			$z = \ON<6->{\B{t^3 + t - t = t^3}}$
		\end{tabular}
	\end{column}
	\begin{column}{0.3cm}
		$\D{\implies}$
	\end{column}
	\begin{column}{6.3cm}
		${\bf r}(t) = 
		\ON<7->{\B{[\,t^3 + t\,,\,t\,,\,t^3\,]}}$ \qquad $t \in 
		\ON<8->{\B{\RE}}$
	\end{column}
	\end{columns}
	
	\vspace{1.4cm}
		
	\EX $x^2 + y^2 = 16$ , $z = 2x + y$ \quad (Intersection of 
	\ON<9->{\B{\,\, a cylinder and plane \,\,}})
	
	\vspace{0.2cm}
	
	Let \,\ON<10->{\B{$t$ be the polar angle in the $xy$-plane:}} \,
	
	\vspace{0.2cm}
	
	\begin{columns}
	\begin{column}{1.0cm}
		%--------TABLE---------
		\begin{tabular}{l}
			$x = \ON<11->{\B{4\cos t}}$ \\
			$y = \ON<12->{\B{4\sin t}}$ \\
			$z = \ON<13->{\B{8\cos t + 4\sin t}}$
		\end{tabular}
	\end{column}
	\begin{column}{0.3cm}
		$\D{\implies}$
	\end{column}
	\begin{column}{6.7cm}
		${\bf r}(t) = 
		\ON<14->{\B{[\,4\cos t\,,\,4\sin t\,,\,8\cos t + 4\sin t\,]}}$ \qquad $t \in 
		\ON<15->{\B{[0,2\pi]}}$
	\end{column}
	\end{columns}
	
	\vspace{3.9cm}
	
\end{column}
\begin{column}{0.0cm}\end{column}
\end{columns}
\end{frame}

%~~~~~~~~~~~~~~~~~~~~~~~~~~~~~~~~~~~~~~~~~~~~~~~~~~~~~~~~~
\begin{frame}
{}
\scriptsize
\begin{columns}
\begin{column}{12.0cm}
\EX Find the velocity function given the position function \quad $\RT = [
\only<3->{\US{\B{ x(t) }}}{
\only<3->{\UB}{ 1+t^3 }}, 
\only<5->{\US{\B{ y(t) }}}{
\only<5->{\UB}{ te^{-t} }}, 
\only<7->{\US{\B{ z(t) }}}{
\only<7->{\UB}{ \sin2t }}]$

\vspace{0.2cm}

Recall: $velocity = \RPT = 
\ON<2->{\B{[ \,x'(t),\,}}
\ON<4->{\B{y'(t),\,}}
\ON<6->{\B{z'(t)\,]}}$
\vspace{0.6cm}

\vspace{0.2cm}
	
	\begin{columns}
	\begin{column}{1.0cm}
		%--------TABLE---------
		\begin{tabular}{l}
			$\ON<7->{\B{x'(t) = 3t^2}}$ \\
			$\ON<8->{\B{y'(t) = e^{-t} - te^{-t}}}$ \\
			$\ON<9->{\B{z'(t) = 2\cos(2t)}}$
		\end{tabular}
	\end{column}
	\begin{column}{0.3cm}
		\ON<10->{\B{$\D{\implies}$}}
	\end{column}
	\begin{column}{6.3cm}
		$ 
		\ON<11->{\B{\RPT = [\,3t^2\,,\,e^{-t} - te^{-t}\,,\,2\cos(2t)\,]}}$
	\end{column}
	\end{columns}
	
	\vspace{1.0cm}
	
\EX Let $\RT = [\sqrt{t}, 2-t]$.  \quad Find ${\bf r}(1)$ and ${\bf r}'(1)$

\vspace{0.2cm}

\begin{columns}
\begin{column}{6.0cm}
	
	$\D{
	\ON<12->{\B{{\bf r}(1) = [\sqrt{1}, 2-1]}}
	\ON<13->{\B{ = [1,1]}}}$

	\vspace{0.7cm}

	$\D{
	\ON<14->{\B{\RPT = \left[\frac{1}{2}t^{-1/2}, -1\right]}}}$

	\vspace{0.7cm}

	$\D{
	\ON<15->{\B{{\bf r}'(1) = \left[\frac{1}{2}(1)^{-1/2}, -1\right]}}
	\ON<16->{\B{ = \left[\frac{1}{2},-1\right]}}}$

\end{column}
\begin{column}{3.0cm}
	\ON<17->{\includegraphics[scale = 0.2]{Figures/9_5_Fig2.png}}
\end{column}
\end{columns}

\end{column}
\begin{column}{0.0cm}\end{column}
\end{columns}
\end{frame}

%%~~~~~~~~~~~~~~~~~~~~~~~~~~~~~~~~~~~~~~~~~~~~~~~~~~~~~~~~~
%\begin{frame}
%{}
%\scriptsize
%\begin{columns}
%\begin{column}{12.0cm}
%	\DEF Line Integral
%\end{column}
%\begin{column}{0.0cm}\end{column}
%\end{columns}
%\end{frame}
%
%%~~~~~~~~~~~~~~~~~~~~~~~~~~~~~~~~~~~~~~~~~~~~~~~~~~~~~~~~~
%\begin{frame}
%{}
%\scriptsize
%\begin{columns}
%\begin{column}{12.0cm}
%	Back to Work
%\end{column}
%\begin{column}{0.0cm}\end{column}
%\end{columns}
%\end{frame}
%
%%~~~~~~~~~~~~~~~~~~~~~~~~~~~~~~~~~~~~~~~~~~~~~~~~~~~~~~~~~
%\begin{frame}
%{}
%\scriptsize
%\begin{columns}
%\begin{column}{12.0cm}
%	
%\end{column}
%\begin{column}{0.0cm}\end{column}
%\end{columns}
%\end{frame}

%%~~~~~~~~~~~~~~~~~~~~~~~~~~~~~~~~~~~~~~~~~~~~~~~~~~~~~~~~~
%\begin{frame}
%{}
%\scriptsize
%\begin{columns}
%\begin{column}{12.0cm}
%{\bf Parameterizing a straight line curves}.  This is simple if you know: 
%
%\vspace{0.2cm}
%
%(1) \ON<2->{\B{a vector, {\bf v}, parallel to the line and}} 
%
%\vspace{0.1cm}
%
%(2) \ON<3->{\B{a point, P, on the line.}}
%
%\vspace{0.4cm}
%
%Then the parameterization is given by:
%
%\vspace{0.2cm}
%
%\begin{center}
%	$\D{\ON<4->{\B{\RT = {\bf v} t + P}}}$ \quad where \quad 
%	\ON<5->{\B{$t \in \RE$}}
%\end{center}
%
%\EX Suppose we know that a line is parallel to $[1,-1,0]$ and passes through $(0,5,5)$.  Then the parameterization for this line is:
%
%\begin{center}
%	$\D{\RT = \, 
%	\ON<6->{\B{[1,-1,0] t + (0,5,5)}}}$ \quad or \quad 
%	\ON<7->{\B{$\D{\RT = \, [t,-t+5,5]}$ \quad $t \in \RE$}}
%\end{center}
%
%\vspace{0.2cm}
%
%\EX Suppose we instead know that a line passes through the points $(-3,1,9)$ and $(-2,2,5)$.  Then we first find the \ON<8->{\B{vector parallel to the line by:}} 
%
%\begin{center}
%	\ON<9->{\B{$\D{{\bf v} = \, (-3,1,9)-(-2,2,5) = [-1,-1,4]}$}}
%\end{center}
%
%and so the parameterization is
%
%\begin{center}
%	$\D{\RT = \, 
%	\ON<10->{\B{[-1,-1,4] t + (-3,1,9)}}}$ \quad or \quad 
%	\ON<11->{\B{$\D{\RT = \, [-t-3,-t+1,4t+9]}$ \quad $t \in \RE$}}
%\end{center}
%
%\end{column}
%\begin{column}{0.0cm}\end{column}
%\end{columns}
%\end{frame}

%~~~~~~~~~~~~~~~~~~~~~~~~~~~~~~~~~~~~~~~~~~~~~~~~~~~~~~~~~
\begin{frame}
{}
\scriptsize
\begin{columns}
\begin{column}{12.0cm}
{\bf Parameterizing a straight line segments}.  In this case we need just the initial and end points of the segment, say $P$ and $Q$.

\vspace{0.4cm}

Then the parameterization is given by:

\vspace{0.2cm}

\begin{center}
	$\D{\RT = 
	\ON<2->{\B{(Q-P) t + P}}}$ \quad where \quad 
	\ON<3->{\B{$t \in [0,1]$}}
\end{center}

\vspace{0.2cm}

\EX Suppose we wanted to parameterize a line with initial point $(-1,-1,-9)$ and end point $(10,0,15)$.  Then the parameterization is given by: 

%--------EQN----------
\begin{align*}
	{\bf v} 
	&= \, \ON<4->{\B{[(10,0,15)-(-1,-1,-9)]}} \, 
	\ON<5->{\B{t}} 
	\ON<6->{\B{\, + \, (-1,-1,-9)}} \hspace{1.0cm} \text{where} \, t \in [0,1] \\
	&\ON<7->{\B{= \, [11,1,24] \, t + (-1,-1,-9)}} \\
	&\ON<8->{\B{= \, [11t-1,t-1,24-9t]}} \\
\end{align*}
%----------------------

\vspace{2.2cm}

\end{column}
\begin{column}{0.0cm}\end{column}
\end{columns}
\end{frame}


%~~~~~~~~~~~~~~~~~~~~~~~~~~~~~~~~~~~~~~~~~~~~~~~~~~~~~~~~~
\begin{frame}
\begin{columns}
\begin{column}{12.0cm}

{II. Review: Integration}

\vspace{0.2cm}

\scriptsize

Remember, integration is another way of saying ``add up a lot of small areas to get a whole area''. So lets think about it as such: \\

\vspace{0.2cm}

\begin{itemize}
	\item Let $dA$ be a small area (of a rectangle) and $A$ the whole area.
	\vspace{0.2cm}
	\item then \quad $\D{dA = \ON<2->{\B{f(x) \cdot dx \quad \text{(see figure) \quad }}}}$ and so
	\vspace{0.2cm}
	\item \quad \,A = \, \ON<3->{\B{$\D{\int_{a}^{b} dA}$}} \ON<3->{\B{$= \D{\int_{a}^{b} f(x) \,dx}$ }}
\end{itemize}

\begin{center}
	\includegraphics[scale = 0.23]{Figures/7_1_Fig2.png}
\end{center}
	
Note: The limits of integrations ($a$ \& $b$) tell you \ON<4->{\B{the range of the small pieces (in this example, $dx$)}}
\end{column}
\begin{column}{0.0cm}\end{column}
\end{columns}
\end{frame}

%%~~~~~~~~~~~~~~~~~~~~~~~~~~~~~~~~~~~~~~~~~~~~~~~~~~~~~~~~~
%\begin{frame}
%{}
%\scriptsize
%\begin{columns}
%\begin{column}{12.0cm}
%
%In the case of area between two curves:
%
%\vspace{0.2cm}
%
%\begin{itemize}
%	\item \quad $\D{dA = \ON<2->{\B{\overset{\text{Area of Rectangle}}{\overset{\text{to $f$ curve}}{\overbrace{f(x) \cdot dx}}} - \overset{\text{Area of Rectangle}}{\overset{\text{to $g$ curve}}{\overbrace{g(x) \cdot dx}}} = \left[f(x) - g(x)\right]\, dx \quad \text{(see figure) \quad }}}}$ and so
%	\vspace{0.2cm}
%	\item \quad \,A = \, \ON<3->{\B{$\D{\int_{a}^{b} dA}$}} \ON<3->{\B{$= \D{\int_{a}^{b} \left[f(x) - g(x)\right]\, dx}$ }}
%\end{itemize}
%
%\vspace{0.4cm}
%
%\centering
%\includegraphics[scale = 0.24]{Figures/7_1_Fig3.png}\\
%	
%\end{column}
%\begin{column}{0.0cm}\end{column}
%\end{columns}
%\end{frame}


%%~~~~~~~~~~~~~~~~~~~~~~~~~~~~~~~~~~~~~~~~~~~~~~~~~~~~~~~~~
%\begin{frame}
%{}
%\scriptsize
%\begin{columns}
%\begin{column}{12.0cm}
%
%Suppose we wanted the area under $\sin x$ for $x$ between $0$ and $\pi$ \\
%
%\vspace{0.2cm}
%
%\begin{center}
%\only<1>{\includegraphics[scale = 0.25]{Figures/7_1_Fig4.png}}
%\only<2>{\includegraphics[scale = 0.25]{Figures/7_1_Fig5.png}}
%\only<3->{\includegraphics[scale = 0.25]{Figures/7_1_Fig6.png}}
%\end{center}
%$$
%dA = \ON<4->{\B{\overset{height}{\overbrace{\sin x}} \, \overset{width}{\overbrace{dx}}}}
%$$
%$$
%A = \int \limits_{0}^{\pi} dA = \ON<5->{\B{\int \limits_{0}^{\pi} \sin x \, dx}}
%$$
%	
%\end{column}
%\begin{column}{0.0cm}\end{column}
%\end{columns}
%\end{frame}

%%~~~~~~~~~~~~~~~~~~~~~~~~~~~~~~~~~~~~~~~~~~~~~~~~~~~~~~~~~
%\begin{frame}
%{}
%\scriptsize
%\begin{columns}
%\begin{column}{12.0cm}
%
%Now, suppose we wanted the area between $\sin x$ and $\frac{1}{2}\sin x$ over $[0,\pi]$ \\
%
%\vspace{0.2cm}
%
%\begin{center}
%\only<1>{\includegraphics[scale = 0.25]{Figures/7_1_Fig7.png}}
%\only<2>{\includegraphics[scale = 0.25]{Figures/7_1_Fig10.png}}
%\only<3>{\includegraphics[scale = 0.25]{Figures/7_1_Fig8.png}}
%\only<4->{\includegraphics[scale = 0.25]{Figures/7_1_Fig9.png}}
%\end{center}
%$$
%dA = (\,\,\ON<3->{\B{\overset{taller}{\overbrace{\sin x}}}} \ON<4->{\B{\,- \overset{shorter}{\overbrace{\frac{1}{2}\sin x}}}}\,\,)\, dx 
%$$
%$$
%Area = \int \limits_{0}^{\pi} dA = \ON<5->{\B{\int \limits_{0}^{\pi} (\sin x - \frac{1}{2}\sin x)\, dx}}
%$$
%	
%\end{column}
%\begin{column}{0.0cm}\end{column}
%\end{columns}
%\end{frame}

%%~~~~~~~~~~~~~~~~~~~~~~~~~~~~~~~~~~~~~~~~~~~~~~~~~~~~~~~~~
%\begin{frame}
%\scriptsize
%\begin{columns}
%\begin{column}{12.0cm}
%{\bf Arc Length Formula:} \\
%\vspace{0.2cm}
%Add up a collection of straight line segments to approximate the total length of the curved line.  \\
%
%\vspace{0.2cm}
%In calculus, we integrate an arbitrary small segment from $a$ to $b$ to get the length of the whole curve.   \\
%
%\includegraphics[scale = 0.22]{Figures/7_4_Fig7.png} \hspace{0.5cm}\includegraphics[scale = 0.20]{Figures/7_4_Fig6.png}	\\
%
%\vspace{0.2cm}
%
%$$
%length = s = \int \limits_{a}^{b} \, ds
%$$
%\end{column}
%\begin{column}{0.0cm}\end{column}
%\end{columns}
%\end{frame}

%%~~~~~~~~~~~~~~~~~~~~~~~~~~~~~~~~~~~~~~~~~~~~~~~~~~~~~~~~~
%\begin{frame}
%\scriptsize
%\begin{columns}
%\begin{column}{12.0cm}
%
%\begin{columns}
%\begin{column}{6.0cm}
%	\includegraphics[scale = 0.2]{Figures/7_4_Fig5.png} \quad $\D{f'(x)=\frac{dy}{dx}}$\\
%	\vspace{0.4cm}
%	$\D{ds = \sqrt{dx^2 + dy^2} = \sqrt{dx^2\left( 1 + \frac{dy^2}{dx^2} \right)}}$ \\
%	$\D{ \hspace{2.18cm} = \sqrt{1 + \left[\frac{dy}{dx}\right]^2}\sqrt{dx^2} }$\\
%	$\D{ \hspace{2.18cm} = \sqrt{1 + \left[f'(x)\right]^2}\,\,dx }$
%	\vspace{0.2cm}
%\end{column}
%\begin{column}{5.0cm}
%	\includegraphics[scale = 0.2]{Figures/7_4_Fig8.png}\\
%	\vspace{1.2cm}
%\end{column}
%\end{columns}
%
%\vspace{0.5cm}
%
%So the arc length of the curve on $f$ (as a function of $x$) from $a$ to $b$ is
%$$
%s = \int \limits_{a}^{b} \, ds = \int \limits_{a}^{b} \sqrt{1 + [f'(x)]^2}\, dx
%$$
%
%	
%\end{column}
%\begin{column}{0.0cm}\end{column}
%\end{columns}
%\end{frame}

%%%%~~~~~~~~~~~~~~~~~~~~~~~~~~~~~~~~~~~~~~~~~~~~~~~~~~~~~~~~~
%%%\begin{frame}
%%%{}
%%%\scriptsize
%%%\begin{columns}
%%%\begin{column}{12.0cm}
%%%
%%%{\bf Volume of rotation formulas} \\
%%%\vspace{0.1cm}
%%%
%%%\begin{columns}
%%%\begin{column}{4.6cm}
%%%1. The disk method \\
%%%
%%%\vspace{0.1cm}
%%%\qquad $\circ$\,\, disk-shaped cross-sections \\
%%%\vspace{1.0cm}
%%%
%%%Adding up small volumes:
%%%\end{column}
%%%\begin{column}{3.0cm}
%%%	\includegraphics[scale = 0.12]{Figures/7_2_Fig14.png}
%%%\end{column}
%%%\end{columns}
%%%
%%%\vspace{0.2cm}
%%%
%%% \quad $Volume \, of \, a \, small \, disk = (
%%% \ON<2->{\B{area \, of \, face}})(
%%% \ON<3->{\B{small \, thickness}})$ 
%%% 
%%% \vspace{0.2cm}
%%% 
%%% \begin{columns}
%%%\begin{column}{3.0cm}
%%%	
%%%	\ON<4->{\B{$area \, of \, face = $}}
%%%	\ON<5->{\B{$\pi R(x)^2$}}
%%% 	
%%%	\vspace{0.2cm}
%%%	
%%% 	\ON<6->{\B{$small \, thickness = $}}
%%% 	\ON<7->{\B{$dx$}}
%%% 
%%%\end{column}
%%%\begin{column}{5.0cm} 
%%%	\ON<8->{\B{$\implies \qquad \D{dV = \pi R(x)^2 dx}$}}
%%%\end{column}
%%%\end{columns}
%%% 
%%% \vspace{0.4cm}
%%% 
%%%\begin{center}
%%%	$\D{Volume = \int \limits_{a}^{b} \, dV = 
%%%	\ON<9->{\B{\int \limits_{a}^{b} \, \pi \left(R(x)\right)^2 dx}}}$ 
%%%\end{center}
%%%
%%%\vspace{3.0cm}
%%%
%%%\end{column}
%%%\begin{column}{0.0cm}\end{column}
%%%\end{columns}
%%%\end{frame}
%%%
%%%%~~~~~~~~~~~~~~~~~~~~~~~~~~~~~~~~~~~~~~~~~~~~~~~~~~~~~~~~~
%%%\begin{frame}
%%%{}
%%%\scriptsize
%%%\begin{columns}
%%%\begin{column}{12.0cm}
%%%
%%%%%
%%%\begin{columns}
%%%\begin{column}{4.6cm}
%%%2. The washer method \\
%%%\vspace{0.1cm}
%%%\qquad $\circ$\,\, washer-shaped cross-sections \\
%%%\vspace{1.0cm}
%%%
%%%Adding up small volumes:
%%%\end{column}
%%%\begin{column}{3.0cm}
%%%	\includegraphics[scale = 0.12]{Figures/7_2_Fig23.png}
%%%\end{column}
%%%\end{columns}
%%%%%
%%%
%%% \quad $Volume \, of \, a \, small \, washer = (
%%% \ON<2->{\B{area \, of \, face}})(
%%% \ON<3->{\B{small \, thickness}})$ 
%%%  
%%% \vspace{0.1cm}
%%% 
%%% \begin{columns}
%%% \begin{column}{0.5cm}
%%% \end{column}
%%%\begin{column}{6.7cm}
%%%	
%%%	\ON<4->{\B{$area \, of \, face = $}}
%%%	\ON<5->{\B{$area \, of \, disk - area \, of \, hole = $}}
%%%	\ON<6->{\B{$\pi R(x)^2 - \pi r(x)^2$}} \\
%%%	$\, \hspace{4.33cm}$ 
%%%	\ON<7->{\B{$ = \pi \left(R(x)^2 - r(x)^2\right)$}}
%%% 	
%%%	\vspace{0.1cm}
%%%	
%%% 	\ON<8->{\B{$small \, thickness = $}}
%%% 	\ON<9->{\B{$dx$}}
%%% 
%%%\end{column}
%%%\begin{column}{4.4cm} 
%%%	\ON<10->{\B{$\implies \,\, \D{dV = \pi \left(R(x)^2 - r(x)^2\right) dx}$}}
%%%\end{column}
%%%\end{columns}
%%% 
%%% \vspace{0.8cm}
%%% 
%%%\begin{center}
%%%	$\D{Volume = \int \limits_{a}^{b} \, dV = 
%%%	\ON<11->{\B{\int \limits_{a}^{b} \, \pi \left(R(x)^2 - r(x)^2\right) dx}}}$ 
%%%\end{center}
%%%
%%%\vspace{3.0cm}
%%%
%%%\end{column}
%%%\begin{column}{0.0cm}\end{column}
%%%\end{columns}
%%%\end{frame}
%%%
%%%%~~~~~~~~~~~~~~~~~~~~~~~~~~~~~~~~~~~~~~~~~~~~~~~~~~~~~~~~~
%%%\begin{frame}
%%%{}
%%%\scriptsize
%%%\begin{columns}
%%%\begin{column}{12.0cm}
%%%
%%%%%
%%%\begin{columns}
%%%\begin{column}{6.2cm}
%%%3. The shell method \\
%%%\vspace{0.1cm}
%%%\quad $\circ$\,\, tube-shaped cross-sections \\
%%%
%%%\vspace{0.3cm}
%%%
%%%\quad $-$\,\, p(x) = distance from shell to axis of rotation. \\
%%%
%%%\quad $-$\,\, h(x) = height of shell. \\
%%%
%%%\vspace{0.2cm}
%%%
%%%
%%%Adding up small volumes:
%%%\end{column}
%%%\begin{column}{2.0cm}
%%%	\includegraphics[scale = 0.18]{Figures/7_2_Fig28.png}
%%%\end{column}
%%%\end{columns}
%%%%%
%%%
%%%\vspace{0.2cm}
%%%
%%% \quad $Volume \, of \, a \, small \, shell = (
%%% \ON<2->{\B{height}})(
%%% \ON<3->{\B{circumference}})(
%%% \ON<4->{\B{small \, thickness}})$
%%% 
%%%\vspace{0.2cm}
%%% 
%%% \begin{columns}
%%%\begin{column}{3.0cm}
%%%	
%%%	\ON<5->{\B{$height = $}}
%%%	\ON<6->{\B{$h(x)$}}
%%% 	
%%%	\vspace{0.2cm}
%%%	
%%%	\ON<7->{\B{$circumference = $}}
%%%	\ON<8->{\B{$2\pi p(x)$}}
%%% 	
%%%	\vspace{0.2cm}
%%%	
%%% 	\ON<9->{\B{$small \, thickness = $}}
%%% 	\ON<10->{\B{$dx$}}
%%% 
%%%\end{column}
%%%\begin{column}{5.0cm} 
%%%	\ON<11->{\B{$\implies \qquad \D{dV = 2\pi \, p(x) \, h(x) \, dx}$}}
%%%\end{column}
%%%\end{columns}
%%% 
%%% \vspace{0.4cm}
%%% 
%%%\begin{center}
%%%	$\D{Volume = \int \limits_{a}^{b} \, dV = 
%%%	\ON<12->{\B{\int \limits_{a}^{b} \, 2\pi \, p(x) \, h(x) \, dx}}}$ 
%%%\end{center}
%%%
%%%\vspace{3.0cm}
%%%
%%%\end{column}
%%%\begin{column}{0.0cm}\end{column}
%%%\end{columns}
%%%\end{frame}

%~~~~~~~~~~~~~~~~~~~~~~~~~~~~~~~~~~~~~~~~~~~~~~~~~~~~~~~~~
\begin{frame}
{}
\scriptsize
\begin{columns}
\begin{column}{12.0cm}
{\bf Interpreting Integrals}: Identify the physical meaning of the Integral

\vspace{0.2cm}

\OBUL $\D{\iiint \limits_{E}^{} \rho(x,y,z)\, dV}$  \quad where \quad $E$ is a solid sphere and $\rho(x,y,z)$ is the density at any point in $E$

\vspace{0.2cm}

\begin{center}
	\ON<2->{\B{Total Mass of E}}
\end{center}

\vspace{0.6cm}

\OBUL $\D{\int \limits_{a}^{b} f(x)\, dx}$  \quad where \quad $f(x)$ is the force required to pull a wagon from $a$ to $b$

\vspace{0.2cm}

\begin{center}
	\ON<3->{\B{Total Work required to pull the wagon a distance $b-a$}}
\end{center}

\vspace{0.6cm}

\OBUL $\D{\int \limits_{C} |{\bf v}(t)|\, dt}$  \quad where \quad ${\bf v}(t)$ is the velocity of a particle traveling on a finite curve $C$

\vspace{0.2cm}

\begin{center}
	\ON<4->{\B{Total distance traveled by the particle along $C$ (i.e. the \ON<5->{\B{\,\,\, arclength\,\,\,}} of $C$)}}
\end{center}

\vspace{0.6cm}

\end{column}
\begin{column}{0.0cm}\end{column}
\end{columns}
\end{frame}

%~~~~~~~~~~~~~~~~~~~~~~~~~~~~~~~~~~~~~~~~~~~~~~~~~~~~~~~~~
\begin{frame}
\begin{columns}
\begin{column}{12.0cm}

{III. Arclength of a Curve}

\vspace{0.2cm}

\scriptsize

\begin{columns}
\begin{column}{4.0cm}

\includegraphics[scale = 0.2]{Figures/9_5_Fig7.png}

\vspace{0.5cm}

\includegraphics[scale = 0.2]{Figures/9_5_Fig5.png}
	
\end{column}
\begin{column}{6.0cm}

Derive the integral expression for $s(t)$.  

\vspace{0.2cm}

Assume $C$ is parametrized by ${\bf r}(t)$ \quad for $a\le t \le b$

\vspace{0.2cm}

We want to think of adding up many ``small" line segments,  so that we can use an integral

\begin{center}
	$\D{s = \int \limits_{C}^{} \, ds}$
\end{center}

But, what is $ds$? \,\, \ON<2->{\B{($ds$ is a ``small" distance)}}

\vspace{0.2cm}

Recall: distance = (rate)(time), so 

\begin{center}
	\ON<3->{\B{``small" distance = \,}} \ON<4->{\textcolor{red}{(rate)}} \ON<5->{\textcolor{orange}{(``small" time)}}
	
	\vspace{0.2cm}
	
	or \quad $\D{\ON<6->{\B{ds = \,}}\ON<7->{\textcolor{red}{|{\bf r}'(t)|}}\,\ON<8->{\textcolor{orange}{dt}}}$
\end{center}
Therefore
\begin{center}
	$\D{s(t) = \ON<9->{\B{\int \limits_{a}^{b} |{\bf r}'(t)|\, dt}}}$
\end{center}

\end{column}
\end{columns}

\end{column}
\begin{column}{0.0cm}\end{column}
\end{columns}
\end{frame}

%~~~~~~~~~~~~~~~~~~~~~~~~~~~~~~~~~~~~~~~~~~~~~~~~~~~~~~~~~
\begin{frame}
{}
\scriptsize
\begin{columns}
\begin{column}{12.0cm}

Compute the arclength of the curve: \quad ${\bf r}(t) = [\cos t, \sin t, t]$ \quad %for \quad $t \in [0,2\pi]$

\begin{center}
	\includegraphics[scale = 0.2]{Figures/9_5_Fig6.png}	
\end{center}

\vspace{0.2cm}

Need: 
\ON<2->{\B{\RPT \, and \, [a,b] \quad (Range of the parameter $t$) }}

\vspace{0.2cm}

Note: 
\ON<3->{\B{Since $z=t$, then the range of $t$-values is the same as the range of $z$-values}} 

\vspace{0.2cm}

\ON<4->{\B{$\implies \quad t \in [0,2\pi]$  \quad (This is $[a,b]$)}}

\vspace{0.2cm}

$\D{
\ON<5->{\B{\RPT = [
\ON<6->{\B{\,\,-\sin t\,\,}}, 
\ON<7->{\B{\,\,\cos t\,\,}} , 
\ON<8->{\B{\,\,1\,\,}}]}}
}$

\vspace{0.2cm}

$\D{
\ON<9->{\B{||\RPT||^2 = \,}} 
\ON<10->{\B{\sin^2 t + \cos^2 t + 1^2 =\,}} 
\ON<11->{\B{1+1 = 2}} \quad 
\ON<12->{\B{\implies \quad ||\RPT|| = \sqrt{2}}}}$

\begin{center}
	$\D{
	\ON<13->{\B{s(t) = \int \limits_{a}^{b} |{\bf r}'(t)|\, dt = \,}} 
	\ON<14->{\B{\int \limits_{0}^{2\pi} \sqrt{2}\, dt = \,}}
	\ON<15->{\B{\left. \sqrt{2} \,t \,\right|_0^{2\pi} = 2\sqrt{2}\pi}}}$
\end{center}

\vspace{4.2cm}

\end{column}
\begin{column}{0.0cm}\end{column}
\end{columns}
\end{frame}

%%~~~~~~~~~~~~~~~~~~~~~~~~~~~~~~~~~~~~~~~~~~~~~~~~~~~~~~~~~
%\begin{frame}
%{}
%\scriptsize
%\begin{columns}
%\begin{column}{12.0cm}
%Additional Homework Problem for section 9.5:
%
%\vspace{0.2cm}
%
%Dervive an integral expression which represents the volume of a pyramid with a square base.  Let the height be $h$ and the side of the base be $b$.
%
%\vspace{0.5cm}
%
%\includegraphics[scale = 0.25]{Figures/9_5_Fig8.png}
%
%\vspace{3.0cm}
% 	
%\end{column}
%\begin{column}{0.0cm}\end{column}
%\end{columns}
%\end{frame}

%%~~~~~~~~~~~~~~~~~~~~~~~~~~~~~~~~~~~~~~~~~~~~~~~~~~~~~~~~~
%\begin{frame}
%{}
%\scriptsize
%\begin{columns}
%\begin{column}{12.0cm}
%Suppose $C$ is some curve in $\RE^3$ and let ${\bf r}(t)$ represent the parameterization of $C$ for $t$ in $[a,b]$.	
%\end{column}
%\begin{column}{0.0cm}\end{column}
%\end{columns}
%\end{frame}
